%======================================================================
\section{導入}
\label{sec:intro}

本論文は、自律的に振る舞う認知システムが長期にわたり存続するための
構造的条件を、数理的な枠組みのもとで記述することを目的とする。
本論文では、自己と環境を区別し、入力を受けて行動を選択し、
内部状態と内部モデルを更新しながら存続を図る系を、
自律認知システム（ACS：Autonomous Cognitive System）として抽象化する。
人間、自律的AIシステム、および後続節で構造的に適用される
組織的システムは、この抽象化の適用先に含まれうるが、
各領域における実証的な妥当性は本論文の主張に含めない。
本論文は、ACSの存続を支える条件を、
規範的な命令としてではなく、可行性に関する構造的な制約軸として定式化する。

本論文が扱う範囲は Phase~1、
すなわち決定論的・離散時間・コンパクト空間・単一ACSの内部動学を
中心とする枠組みに限定される。
確率的拡張、連続時間化、社会動態の本格的展開は後続の段階に委譲する。
以下ではまず背景と動機を述べ（\S\ref{subsec:background}）、
本論文の立場とスコープ、および構造・機能概念の分類規則を定め
（\S\ref{subsec:scope}）、
主要概念を予告したうえで（\S\ref{subsec:preview}）、
本論文の貢献（\S\ref{subsec:contributions}）、
関連研究（\S\ref{subsec:related}）、
全体構成（\S\ref{subsec:roadmap}）を順に述べる。

%----------------------------------------------------------------------
\subsection{背景と動機}
\label{subsec:background}

社会のなかで生きる主体の行動は、
熟慮による選択として現れることもあれば、
熟慮的な選択過程を経ない局所的・反応的・自動的な行動として
現れることもある。
こうした行動は、多くの場合、
その主体が属する社会や文化の規範に照らして評価される。
しかし規範は社会や文化ごとに異なり、
同じ行動が一方では妥当とされ、他方では非難されることがある。
規範による評価だけに依拠すると、
異なる規範のもとにある主体どうしが、
互いの行動を共有可能な形で理解する手がかりを得にくい。
人間と自律的AIシステムのように、
成り立ちの異なる主体が同じ社会のなかで行動する場合には、
この問題はいっそう明確になる。

本論文の関心は、
主体の行動がどのような構造的条件と結びついて成立するのかを、
特定の社会や文化の善悪の判断を前提としない形式で記述することにある。
本論文では、行動の成立に関わる可行性上の要請、
内部状態と内部モデル、外部環境、および内部機構の作動の組を、
その行動の構造的源泉と呼ぶ。
構造的源泉は、行動の単一の原因を指す語ではない。
この記述のために本論文は、
人間、自律的AIシステムその他の自律的に振る舞う系を、
共通の抽象構造をもつ自律認知システムとして定義する。
この共通構造のもとでは、
主体の行動を善悪の基準によって判定するのではなく、
どの可行性上の要請が関わり、
どの内部機構の作動を経てその行動が成立したかとして記述できる。
熟慮的な評価を経る応答と、
局所的な保護応答や自動化された処理による応答も、
同じ構造のもとで区別して記述される（\S\ref{sec:tipa}）。
また、内部モデルや評価過程が関与する行動については、
主体が利用可能な情報にもとづいて形成した評価と、
形式モデル上の可行性とが一致することを要求しない。
このため、誤認や評価の誤差を伴う行動も、
同じ構造的記述の対象となる。
本論文は、このような共有可能な記述言語の
決定論的・構造的な基礎を与えることを目指す。
ここでいう記述は、行動の予測や原因の一意な特定を意味しない。
その射程の限定は \S\ref{subsec:scope} で述べる。

本論文が提示するACS理論
（自律認知システム理論／Autonomous Cognitive System Theory）は、
著者が研究を進める倫理的協調知性理論
（ECIT：Ethical Coordinative Intelligence Theory）の
基本構造理論として位置づけられる。
ECITは、人間とAIを含む複数の主体が、
対立を固定化させずに協調を続けるための枠組みをどう与えるか、
という問いに関わる。
この問いを論じるためには、
協調の当事者となる主体それぞれの行動を、
特定の規範を前提とせずに共有可能な形で記述する基盤が要る。
ACS理論は、その記述の基盤を与える理論として構想されている。
すなわちACS理論は、
三原理である自己維持・環境利用・可行応答、
自己と環境の構造的結合、
および外部対象との関係構造を記述する枠組みをECITに提供する。
これらがECITの協調概念とどのように対応するかは、
Phase~3 で検討する。

このとき本論文が採る基本的着想は、
倫理を善悪の判断としてではなく、
系が可行であり続けるための制約として記述することである。
すなわち、「こうあるべき」という規範的要請を外部から導入するのではなく、
長期にわたり可行性を保つ系がどのような構造を備えるかを記述し、
その記述のなかに倫理の礎となる条件が
構造的に現れるかどうかを問う。
行動の構造的源泉を記述するという本論文の関心は、
この着想と結びついている。
主体の行動を、可行性に関する要請と、
それに応答する内部機構の作動との関係として記述することが、
特定の善悪の判断を前提としない記述言語の中身となるからである。
この立場のもとで、三原理は価値的な命令ではなく、
長期可行性に関わる相異なる構造的制約軸として扱われる。

ECITはACS理論の上位に位置づけられる枠組みであるが、
その本格的な展開は後続の段階の課題とし、
ACS理論との接続は Phase~3 に位置づける。
本論文では、ECITの概要を動機づけの水準で非形式的に述べるにとどめ、
その規範的・社会適用的側面を Phase~1 の主張に混入させない。
ECITの具体的な要素やそれらの機構への対応づけ、
社会的応用に関する主張は、
本論文の形式的帰結としてではなく、
Phase~3 における今後の研究方向として位置づけられる
（\S\ref{subsec:phase_outlook}）。

%----------------------------------------------------------------------
\subsection{本論文の立場とスコープ}
\label{subsec:scope}

本論文は、倫理的判断そのものを与えるものではない。
本論文が扱うのは、ACSが長期にわたり可行性を保つうえで、
どのような構造的条件が要請されるかである。
この意味において、本論文は記述理論であり、規範理論ではない。
また構造理論であり、価値理論ではない。
本論文が記述するのは「ACSはいかにあるべきか」ではなく、
「長期にわたり可行であるACSはいかなる構造を備えるか」である。
倫理に相当する条件は、外部から課される規範としてではなく、
可行性保存上の制約として記述される。
ECITが扱う規範的・社会適用的な側面は、
本論文の主張には含めず、後続の段階へ委譲する。

\S\ref{subsec:background} で述べた行動の記述について、
その射程を次のように限定する。
第一に、記述は主体の行動を予測することでも、
ある行動の原因を一意に特定することでもない。
本論文の枠組みは行動選択の一般法則を与えず、
一つの行動に複数の構造的要請と内部機構が関与しうることを認める。
第二に、記述の言語が特定の文化の善悪の判断を前提としないことと、
個別の主体へ適用する際のモデル化が
観察者の選択から独立であることとは区別される。
可行集合や必要環境要因をどのように設定するかには、
モデル化の選択が介在する。
第三に、本論文の共通構造が適用されるのは、
\S\ref{sec:acs_def} の定義と構成要件を満たす系であり、
個別の領域における経験的な当てはまりは
本論文の主張に含めない（\S\ref{sec:limitations_and_outlook}）。

長期可行性に関する本論文の形式的な解析は、
可行集合 $K$ と、必要環境充足に相対する
$E_{\mathrm{req}}$ ロバスト可行性核 $\mathrm{Viab}_f(K)$
（定義~\ref{def:viab_core}）を中心とする。
可行性核に留まる応答を選び続ける軌道についての結論
（命題~\ref{prop:kernel_maintenance}）は、
必要環境要因の充足と行動の選択に関する前件のもとでの条件付きの結果であり、
本論文は、それらの前件が実在のACSにおいて維持されることを主張しない
（備考~\ref{rem:kernel_response_scope}）。
また本論文は、可行集合への所属と、機能的終端または崩壊の不在との同値を
一般的な帰結として置かない
（備考~\ref{rem:principle_time_domain_scope}、備考~\ref{rem:sm_not_collapse}）。

本論文の形式的内容は Phase~1 に限定される。
Phase~1 は決定論的・離散時間・コンパクト空間を前提とし、
単一ACSの内部動学を中心に扱う。
この枠組みのもとでは、確率的な記述や期待値的な評価は用いず、
有限窓における決定論的な集計量のみを用いる。
Phase~1 の前提と各対象の型を変えずに、
規則・閾値・操作的基準を具体化する補助的詳細化は Phase~1.$\alpha$ とし
（\S\ref{subsec:phase1a}）、
確率的拡張・連続時間化・社会動態の定量化は Phase~2 の課題とし、
ECITとの本格的な接続は Phase~3 の課題とする。

なお、本論文で扱う組織ACSは、
社会動態の一般理論ではなく、
組織自身がACSの定義を満たす場合に、
単一ACSの枠組みを社会的組織へ構造的に適用する記述として位置づけられる。
したがって、社会的相互作用の確率的変動、制度形成、
集団内の長期的進化、ネットワーク構造の動的変化を
一般的に記述することは、Phase~1 の主張には含めない。
社会は、ACSその他の主体が帰属・所属・活動する場として扱う
（定義~\ref{def:org_terms}）。

本論文の各主張は、
明示された仮定のもとで（under stated assumptions）、
本枠組みの範囲内で（within present framework）
成立するものとして述べられる。

\subsubsection{本論文が対象とするACSと「必要」の意味}
\label{subsubsec:real_world_scope}

ACSは、物理的制約のもとで作動する実在のシステムを記述するための
概念である。
生物学的に実現されるACSはもとより、
計算機上で実現されるACSも、
有限の計算資源、有限の応答時間、有限の記憶容量、
供給とエネルギーの制約、
および構成要素の損傷・劣化可能性から自由ではない。
本論文は、この物理的制約の被拘束性をACSの副次的性質ではなく、
その内部機構を論じる際の前提として扱う。
したがって本論文は、物理的に実装されない仮想的なACSを
別の型として立てず、ACSを一律に扱う。

本論文が想定する環境は、
時間変化、部分観測性、有限資源、有限応答時間、
損傷・劣化可能性を伴い、
他主体との相互作用を含みうる環境である。
ここで、他主体との相互作用を「含みうる」とすることは、
Phase~1 において多主体動学を一般的に形式化することを意味しない。
単一ACSから見た他主体の作用は、
必要に応じて外部環境入力の一部として取り扱われ、
多ACS間相互作用そのものの動学的展開は後続段階へ委譲する。

一方、完全に固定された環境、完全情報が与えられる環境、
実質的に無制約な資源をもつ環境その他の退化した設定は、
本論文の考察対象に含めない。
この除外の一部は、後続節において形式的な仮定として与えられる。
すなわち、環境列が時間的に一定でないことは
\S\ref{sec:three_principles} の非定常環境の仮定
（仮定~\ref{ass:nonstationary}）として、
外部環境状態の差異が可行行動集合に一切影響しない退化的な系の除外は
\S\ref{sec:acs_def} の可行行動集合に対する環境入力の非退化性の仮定
（仮定~\ref{ass:env_viable_nondegeneracy}）として、
それぞれ導入される。
上に挙げたその他の退化した設定（情報および資源に関するもの）の除外は、
Phase~1 では考察対象の射程を定める宣言にとどめ、
独立の形式的仮定としては置かない。
なお、現在可行な行動の少なくとも一つについて、
必要環境要因がすべて充足される外部環境状態のうちに
次状態を $K$ に置かないものがあること
（可行な応答の弱点、定義~\ref{def:nonresponse_env_set}）は、
対象クラスを限定する条件ではなく、
\S\ref{sec:three_principles} の応答多様性縮退と長期可行性の条件付き解析において
命題の前提としてだけ用いられる。

以上より、本論文で以下「必要」と記す場合、
別途より強い論理的必然性が証明されている場合を除き、
上記の環境において長期的可行性の維持を指向するACSに対する
構成的必要性を意味する。
上記の退化した設定を用いてこの構成的必要性を否定することは、
本論文の射程外とする。

この射程の限定は、後続節で個別に証明される
定義的必然性、条件付き必然性または数学的必要性の強さを
弱めるものではない。
とくに、\S\ref{sec:three_principles} において
明示的な仮定のもとで証明される
各命題・補題の論理的地位は、
各結果自身の仮定、定義域および限定によって定まり、
本項の射程限定によって置き換えられない。

\subsubsection{構造・機能概念の三分類}
\label{subsubsec:abc_classification}

本論文は、ある機能がACSに要求されることと、
その機能を本論文がどのような形式で表現するかと、
その表現を解析するために導入する理論量とを、
それぞれ別の事柄として扱う。
これらを区別せずに読むと、
本論文が採用する機構分解や内部表現が、
実在ACSに一対一で備わる部位として要求されているかのように
受け取られうる。
そこで次の三分類を導入する。

\begin{definition}[構造・機能概念の三分類]
\label{def:abc_classification}

本論文では、ACSの構造、機能、機構形式および時刻依存量について、
それらの理論上の地位を区別するため、次の三分類を用いる。

\begin{description}

\item[分類A：構成的必要機能]
\S\ref{subsubsec:real_world_scope} で述べた環境において
長期的な可行性の維持を指向するACSに対して、
構成的に要求される機能をいう。
分類Aは機能的必要性を表し、
当該機能が独立した物理部位、計算モジュール、
状態変数または明示的内部表現として
実装されることを要求しない。
同じ機能的要請を満たす機能的に等価な実現を許す。
分類Aへの指定それ自体は、
その必要性があらゆる数学的環境について
証明済みであることを意味しない。
その必要性の射程は \S\ref{subsubsec:real_world_scope} に従う。

\item[分類B：標準機能形式]
分類Aの構成的必要機能その他のACS内部機能を、
本論文の形式体系において分析・比較・記述可能な責務へ分解するために
採用する標準的な機構、機能分解、内部表現、
接続構造または作動形式をいう。
分類Bは単なる例示ではなく、本論文が採用する標準機能形式である。
一方、生物学的または人工的実装において、
各分類B要素が一対一の独立部位または独立モジュールとして
存在することは要求しない。
また、分類Aの機能が機能的に等価な別の実装構成によって
担持される可能性を排除しない。

\item[分類C：記述・診断・比較概念]
ACSの状態、作動水準、時刻間変化、能力、偏り、健全性、
外部対象との関係、および失敗状態への近接その他を
記述、診断、比較または追跡するために
本論文が導入する理論上の構成物をいう。
これには、スカラー量およびベクトル量のほか、
記述子、カテゴリー、状態帯、集合、経路集合その他の
分析上の構成物を含む。
分類Cの概念は、実在ACSが同一の記号、尺度または内部変数として
保持または計算することを要求しない。
分類Cへの指定は、
その概念がACSの7タプルへ新たな構成要素または
独立状態座標を追加することを意味しない。
また本論文は、別途定義・仮定・命題による接続がない限り、
分類Cの特定値、カテゴリーへの所属または状態帯への所属のみから
三原理の時刻別充足述語の真偽、閉塞または崩壊を一般的な帰結として置かない。

\end{description}
\end{definition}

\begin{remark}[三分類の適用規則]
\label{rem:abc_classification_rules}

定義~\ref{def:abc_classification} の三分類は、
以下の規則に従って適用する。

\begin{enumerate}
  \item 分類は、名称の印象ではなく、
        当該の定義・備考等が直接導入する理論対象に対して付与する。

  \item 原則として、一つの理論対象には一つの分類を付与する。
        一つの記述が分類Aの必要機能と分類Bの標準実現形式を
        同時に含む場合には、
        必要機能と実現形式を別の理論対象として分離する。

  \item 分類Aは機能的必要性を表す。
        したがって本論文は、分類Aの機能が必要であることから、
        特定の分類B機構、物理部位または計算モジュールが
        一対一に存在しなければならないことを一般的な帰結として置かない。

  \item 分類Bは本論文が採用する標準機能形式を表す。
        分類Bであることは単なる任意例を意味しないが、
        対応する機能の物理的・生物学的・計算的実装を
        一意に指定するものでもない。

  \item 分類Cは記述・診断・比較のための概念であり、
        ACSの7タプルへ新たな構成要素を追加するものではない。
        実装が分類Cと同型の内部量を実際に保持する場合でも、
        その事実だけで分類を変更しない。

  \item 三分類は、理論対象が本論文内で担う
        機能的・表現的・記述的地位を区別する軸である。
        これに対し、\S\ref{sec:three_principles} で用いる
        定義的必然性、条件付き必然性、構造的要請、
        ならびに個々の定理・命題の証明強度は、
        論理的地位・主張強度を表す別の軸である。
        したがって、分類Aであること自体から
        数学的必要性または定理としての証明済み性は導かれず、
        逆に定理として証明された対象の論理的地位が
        分類A/B/Cによって弱められることもない。
        仮定の導入地位および強さも、三分類とは別の軸である。

  \item 本論文は、分類Cの概念の端点、閾値、特定値、
        カテゴリーへの所属または状態帯・経路集合への所属から、
        別途明示された接続条件なしに
        閉塞 $A_t=\emptyset$ または
        崩壊 $\mathrm{Col}_t=1$ を一般的な帰結として置かない。

  \item 本分類は、主として \S\ref{sec:tipa} 以降で導入される
        機能、機構形式、補助的内部表現、作動形式、
        時刻依存量、関係記述および失敗状態の分析概念に適用する。
        \S\ref{sec:acs_def} のACS 7タプルを構成する基礎集合・写像そのもの、
        \S\ref{sec:three_principles} の三原理、
        および \S\ref{sec:acs_def} における
        閉塞・機能的終端・崩壊そのものの基礎定義は、
        本分類より上位の基礎定義・構造条件または対象現象として扱い、
        原則として分類A/B/Cの表記対象外とする。
        一方、\S\ref{sec:rel_min} の関係記述子・関係カテゴリー等、
        および \S\ref{sec:collapse_cond} の記述的状態帯・経路集合等は、
        その直接の理論的役割に応じて原則として分類Cの対象とする。
\end{enumerate}

後続節では、分類対象の初出となる正式な定義または備考の表題末尾に、
原則として「（分類A）」「（分類B）」「（分類C）」のいずれかを付す。
同一対象を後続箇所で参照するたびに分類表記を反復する必要はない。
\end{remark}

\subsubsection{形式表現とACS内部認識の区別}
\label{subsec:description_levels_global}

本論文で導入する数式、指標、分類、状態帯、経路集合その他の形式対象は、
理論上の記述・分析のための表現である。
本論文は、その存在だけから、
ACS自身が同一の記号・尺度で当該量の真値を内部に保持し、
完全に認識し、または直接計算することを一般的な帰結として置かない。

ACS内部での評価、診断、応答形成または更新に必要な情報と機能は、
各節で明示される機能要件および標準機能形式に従う。
したがって、
論文上の形式表現が整形式に定義されること、
ACS内部で当該対象を測定・認識・利用できること、
および外部観測者が当該対象を測定できることを区別する。

想定する読者は、Viability Theory、凸解析、動的システム論、
Clarke劣微分の基礎を持つ専門家である。
本論文では、これらの基礎的概念については逐一の導入を行わず、
必要な箇所で参照する。

\subsubsection{記号と添字の規約}
\label{subsec:notation}

本論文の記号では、対象を識別する添字を下付きに置き、
種類、個体、成分、窓、時刻の順に並べる
（例えば $d_{\mathrm{eff},i,t}$、$\bar Q_{i\leftarrow j,W_Q,t}$）。
上付きには、演算と、同じ基本記号で表される対象の変種を区別する固定したラベル
（$E^{\mathrm{bg}}$ 等）だけを置く。
括弧内には、写像、対応または演算の入力を置く。
単一個体だけを扱う文脈では個体の添字を省略する。
規約の全体は付録~\ref{app:notation} に示す。

%----------------------------------------------------------------------
\subsection{主要概念の予告}
\label{subsec:preview}

本論文では、ACSを次の7タプルとして定義する。
\[
\Sigma=(X,U,E,f,\Pi,M,\mathrm{Update}).
\]
ここで $X$ は内部状態空間、$U$ は行動空間、
$E$ は必要環境空間、$f$ は状態遷移写像、
$\Pi$ は持続評価関数、$M$ は内部モデル構造、
$\mathrm{Update}$ は内部モデル更新作用素である。
各構成要素の厳密な定義と型は
\S\ref{sec:acs_def} で与える。
持続評価関数 $\Pi$ から派生する
可行集合
\[
K:=\{x\in X\mid\Pi(x)\ge0\}
\]
と、その時刻ごとの可行行動集合は、
本論文で定める可行性に関わるACSの状態と行動を特徴づける中心的な対象となる。
ACSの7タプルおよびその基礎集合・写像は、
定義~\ref{def:abc_classification} の三分類より上位の
基礎定義として扱う。

ACSの長期可行性に関わる制約軸として、
三原理を導入する。
自己維持（SM）は現在状態が可行集合に属すること、
環境利用（EU）は一段階の可行性に関して臨界となる必要環境要因が充足されること、
可行応答（FR）は臨界必要要因が充足される局面で
次状態を可行集合に置く応答が得られることに対応する。
外部環境状態の差異に対する自己と環境の構造的結合は、
三原理のいずれでもなく EU の前提として扱い、
外部対象との関係品質 $Q$ は \S\ref{sec:rel_min} の記述量とする。
これらの厳密な定式化と、
それぞれが長期可行性に対して持つ論理的地位は
\S\ref{sec:three_principles} で扱う。
三原理は、三分類による機能形式の区別より上位の
制約軸として位置づける。

これらの原理水準の条件と、
ACSにおいてそれらに関わる差異を作動へ接続する
具体的な機能構造とを区別するため、
\S\ref{sec:tipa} では、
情報統合機構M1、適応的評価機構M2、応答行動機構M3からなる
TIPA（Triadic Information-Processing Architecture）、
物理的基盤層PSL、
実装部位層ISL（Implementation Substrate Layer）を定式化する。
ここで、本論文が対象とするACSに構成的に要求される機能は
必要に応じて分類Aとして明示し、
TIPA・PSL・ISLにおける具体的な機構分解、
内部表現、接続構造および作動形式は、
それらを本論文の形式体系内で分析するための
分類Bの標準機能形式として区別する。
この区別により、必要機能の存在を、
特定の生物学的部位または人工的モジュールの
一対一の存在と同一視しない。

また、\S\ref{sec:extended_dynamics} では、
\S\ref{sec:tipa} で定義された同一時刻構造を再定義せず、
一段の作動記録が次の時刻へ閉じる標準作動形式を与え、
内部モデル、機能的記憶、偏向性、実効的発展自由度、
予測精度と実効的時間地平、基盤・実装部位の健全度その他について、
時刻 $t$ から $t+1$ への決定論的な変化を記述する。
このうち、作動水準、変化量、偏り、能力、健全性等を
記述・診断・比較するために導入される理論量は、
主として分類Cとして位置づける。
分類Cの量は、実在ACSが同じ数値表現を内部に保持することを
要求するものではない。

以上により本論文では、
「ACSに構成的に必要な機能」、
「本論文が採用する標準機能形式」、
「理論上の記述・診断・比較量」を区別しながら、
同一の形式的枠組みの内部で接続する。

%----------------------------------------------------------------------
\subsection{本論文の貢献}
\label{subsec:contributions}

本論文の貢献は、
自律的システム、知能、倫理、協調を
概念的に関連づけることにあるのではなく、
これらに関わる構造を単一の形式的枠組みの内部で
接続する点にある。
すなわち、可行性の保存、内部機構、
時刻間の内部動学、関係性の記述、
失敗状態の区別を、
決定論的・離散時間の共通語彙のもとで定式化する。
Phase~1 では確率的記述・連続時間動学・
多主体ネットワーク動態を導入せず、
この限定によって理論の最小骨格を明確化する。
具体的な貢献は次の六点である。

第一の貢献は、ACSの形式的定義である。
人間・自律的AIシステム・組織的システムとして
記述されうる対象を共通の形式で扱うために、
ACSを、単なるエージェント、制御器、
あるいは報酬最大化システムとしてではなく、
\S\ref{subsec:preview} に示した7タプルとして定式化した
（\S\ref{sec:acs_def}）。
この定式化のもとで、
自己維持は外部から与えられた目的の達成としてではなく、
可行集合 $K$ とそこから導かれる可行行動集合を通じて表現される。
これにより、
状態が可行集合を外れること、
可行行動が消失すること、
および構造的個体としての同一性を成立させていた組織関係が
同一個体として回復しない形で失われることを
区別して扱うための基礎が与えられる。
また、外部環境状態を外生背景成分と行動応答成分に分け、
行動応答成分が当該ACSの行動に応じて遷移する形を与える
（定義~\ref{def:env_realization}）。
また、内部モデル $M$ とその更新作用素 $\mathrm{Update}$ を
7タプルの内部に明示することで、
ACSを単なる状態遷移系としてではなく、
自己の状態と環境に関する内部的な処理・更新を伴う
自律認知システムとして記述するための形式的基盤を与える。

第二の貢献は、三原理の定式化とその位置づけである。
長期可行性に関わる三つの制約軸として、
自己維持・SM（Self-Maintenance）、
環境利用・EU（Environmental Utilization）、
可行応答・FR（Feasible Response）
を定義した（\S\ref{sec:three_principles}）。
それぞれの時刻別充足述語は、
SMが現在状態の可行性を、
EUが一段階の可行性に関して臨界となる必要環境要因の充足を、
FRが臨界必要要因の充足される局面で次状態を $K$ に置く応答の存在を問う。
SMは可行集合の定義と同値な条件であり、
EUの時刻別充足述語が偽であることは閉塞を伴い、
閉塞は臨界必要要因の充足の有無によって二分される。
外部環境状態の差異に対する構造的環境結合条件は、
三原理のいずれでもなくEUの前提として位置づける。
三原理の時刻別充足述語の間については、
示した条件付き含意と、本論文が主張しない含意とを区別して整理した
（\S\ref{subsec:principle_linkage}）。
あわせて、三原理に関係する状態を有限時間窓で集約する
三原理時間統合評価条件（T-STAB）を定め、
応答多様性の縮退と可行な応答の弱点が持続する場合に頑健応答不能が持続すること、
およびこれに強仮定を加えると有限時間で閉塞に至ることを示した
（\S\ref{subsec:sd_necessity}）。
以上はいずれも価値判断としてではなく、
本枠組み内の構造的条件として扱われる。
多主体系における均衡の扱いは後続の段階の課題として残る。

第三の貢献は、
三原理を抽象的条件に留めず、
本論文が対象とするACSに構成的に要求される機能と、
それらを本論文内で分析可能にする標準機能形式とを
区別したうえで、内部機構の形式的構成を与えた点である
（\S\ref{sec:tipa}）。
本論文は、情報統合、可行性感応評価、応答形成、
基盤維持、入出力接続、実装担体その他の機能を、
必要に応じて分類Aの構成的必要機能として区別する。
そのうえで、これらを分析可能な責務へ分解する
分類Bの標準機能形式として、
情報統合機構M1、適応的評価機構M2、応答行動機構M3からなるTIPA、
物理的基盤層PSL、実装部位層ISLを定式化する。
また、識別、Priority、要請、候補形成、機能的記憶、
高速処理、Automaticityその他の内部表現・作動形式についても、
それらを現実の独立部位と同一視せず、
本論文における標準機能形式として位置づける。
これにより、原理水準の条件と実装可能な内部構造を切り離さず、
かつ特定の生物学的部位や人工的設計へ還元しない形で、
ACSの内部機構を価値中立的に記述する基礎を与える。
さらに、分類A・B・Cの区別を導入することで、
必要機能、標準実現形式、解析上の記述量を
存在論的に同一視しない読み方を明示した。

第四の貢献は、
単一ACSの時刻間の接続を、
同一時刻の機構構造と区別して整理した点である
（\S\ref{sec:extended_dynamics}）。
\S\ref{sec:extended_dynamics} は、
\S\ref{sec:tipa} のTIPA・PSL・ISLの
同一時刻における構造・責務を再定義しない。
原始的な機能の標準写像と読み出し規則を固定し、
環境実現を与えたとき、
一段の作動記録がただ一つ定まり次の時刻へ閉じること、
およびその記録の因果性を示した。
そのうえで、学習と忘却、
偏向性とその程度である偏向強度、
実効的発展自由度、
予測精度係数と実効的時間地平、
PSL・ISLの健全度を、
状態・内部モデル・外部環境状態または作動記録の有限区間からの読み出しとして時刻 $t$ から $t+1$ へ接続する。
同時に、学習による変化と改善、
「内部量の特定値」と「閉塞・崩壊」を
相互に同一視しない。
これにより\S\ref{sec:extended_dynamics}は、
特定の学習則や収束則を仮定するのではなく、
\S\ref{sec:tipa}の同一時刻構造を
後続の関係性・失敗状態解析へ接続する
時刻間のインターフェースを与える。

第五の貢献は、
ACSと外部対象との関係性を、
補助的な社会的記述としてではなく、
関係継続中に記録される可行性関連量の変化を記述する形式として
導入した点である
（\S\ref{sec:rel_min}）。
本論文では関係性の最小構造を定式化し、
観測、依存、支援、協力、共生、支配、管理、
操作、搾取、競争、回避、対立などを
価値中立的なカテゴリーとして整理する。
この関係記述は、外部対象との接続を
単なる有無として扱うのではなく、
関係継続中の有限窓において記録される
可行性、応答可能性、時間的余裕その他の関連量の変化を
記述するための形式的骨格を与える。
これらの差分は、外部対象の単独因果効果を意味しない。
また、評価可能な関係について、
記述者が指定する評価規則に相対する関係品質 $Q$ を導入する。
$Q$ は構造的環境結合条件やEUの条件ではなく、
ACS内部の評価でもない。
一方でPhase~1では、
関係ネットワークの一般動学や
社会的相互作用の確率過程までを主張しない。

第六の貢献は、
失敗状態の区別と組織への構造的適用である。
可行行動が存在しない閉塞
$A_t=\emptyset$
と、
構造的個体としての同一性を成立させていた組織関係が
本枠組みに相対して回復しない形で失われた崩壊
$\mathrm{Col}_t=1$
とを別個の対象として定式化し（\S\ref{sec:acs_def}）、
両者の関係、外的回復可能性、軌道記録と到達経路を整理する
（\S\ref{sec:collapse_cond}）。
これにより、
「現時点で可行な行動が存在しないこと」と、
「同一個体としての構造的自己が回復不能となること」を
同一の失敗概念へ還元しない。
さらに、単一ACSの枠組みを、
単なるACS集団と区別された組織ACSへ適用する
（\S\ref{sec:org_acs}）。
組織自身がACSの定義を満たす場合には、
先行節の閉塞・崩壊・三原理の定義を組織水準へそのまま適用する。
参加する個体ACSの数は組織ACSの成立の条件に含めず、組織ACSの区分の記録として扱う。
構成員変化を含む組織・社会動態の一般化と定量化は Phase~2 の課題とする。

以上の成果の理論的限界と未解決問題は
\S\ref{sec:limitations_and_outlook} で開示され、
Phase~1.$\alpha$（決定論的枠組みのもとでの補助的詳細化）への委譲事項と、
Phase~2（確率的拡張・連続時間化・組織／社会動態）
および Phase~3（ECITとの接続）への展望が整理される。
Phase~1 の目的は、
経験的妥当性や確率的予測を主張することではなく、
倫理を可行性保存制約として記述する立場のもとで、
ACS理論の最小構造を明示し、
後続の段階への接続基盤を与えることにある。
各主張の強さは、
対応する定義・仮定・証明の強さを超えないように制御される。

%----------------------------------------------------------------------
\subsection{関連研究}
\label{subsec:related}

本論文の数理的基盤は、
いくつかの既存の枠組みに依拠する。

可行集合および可行核の概念は、
Viability Theory の枠組み
\cite{Aubin1991,AubinFrankowska1990} に基づく。
ただし、Viability Theory は
可行集合と可行核の一般理論を与えるものであり、
本論文が扱う内部機構の構成、
関係性記述、
閉塞と崩壊の区別といった記述は、
その主題に含まれない。
本論文の位置づけは、
可行性の語彙とこれらの構造記述を
単一の形式的枠組みの内部で接続する点にある。

持続評価関数の劣微分構造は、
非平滑解析における
Clarke の一般化勾配・一般化劣微分の枠組み
\cite{Clarke1990} に依拠して記述される。
ただし、本論文は非平滑解析そのものを拡張するものではなく、
持続評価関数の境界構造を記述するために
既存の数学的道具を用いるにとどまる。

標準的なAIエージェント論
\cite{RussellNorvig2020} は、
知覚に基づいて行動を選択する主体を扱う。
本論文はこの視点と接続するが、
ACSを知覚から行動への写像としてのみ扱うのではなく、
可行性評価、内部モデルの更新、
閉塞、崩壊を備えた構造として扱う点で立場を異にする。

同様に、強化学習 \cite{SuttonBarto2018} は
報酬最大化を中心概念として行動選択を扱うが、
本論文は報酬を第一原理としない。
ACSの存続基準は、
何らかの量の最大化としてではなく、
持続評価関数の符号条件
$\Pi(x_t)\ge0$、
すなわち可行性の保存として与えられる。
報酬に基づく評価・選択は、
TIPAにおける行動比較の一つの実装形態として
読みうるが（\S\ref{sec:tipa}）、
この対応の形式化は本論文の主張には含めない。
また、内部状態の設定点からの距離を駆動とし、
駆動の減少を報酬として扱う恒常性強化学習がある
\cite{KeramatiGutkin2014}。
本論文の自己維持（SM）の時刻別充足述語は、
設定点への接近ではなく閾値に対する符号条件 $\Pi(x_t)\ge0$ であり、
報酬を定めない。

また、自己維持を必須変数の維持として捉える発想は、
Ashby の古典的議論 \cite{Ashby1952} とも接続する。
本論文はこの自己維持の問題意識を共有する一方で、
可行集合、可行行動、内部モデル、
TIPA・PSL・ISL、関係性、
閉塞・崩壊を同一形式体系内で区別して記述することを目指す。

Ashby はさらに、調節の問題を必要多様性の法則として論じた
\cite{Ashby1956}。
\S\ref{sec:three_principles} の応答多様度
（定義~\ref{def:mu_eff}）は、
$E_{\mathrm{req}}(x)$ のうち次状態を $K$ に置く環境の範囲が等しい行動を
一つに数えた同値類の数であり、
外乱の多様性と調節の多様性を比べる法則を述べるものではない。

生命的な自律性の研究では、
系が自己の可行性の条件に対してどのように自身を調節するかが論じられてきた。
Di Paolo は、可行性の条件に関して自己を調節する性質として適応性を論じた
\cite{DiPaolo2005}。
Barandiaran と Egbert は、最小の代謝系のモデルで
可行な領域と不安定な領域を区別し、
可行な領域へ戻すのに必要な環境の変化を用いて、
規範への従属を定義した
\cite{BarandiaranEgbert2014}。
Barandiaran らは、行為主体の定義に
個体性、相互作用の非対称性、規範性の三条件を置いた
\cite{BarandiaranDiPaoloRohde2009}。
これらの三条件は本論文の三原理とは異なる。
本論文は、可行性条件と自律系の行動・調節との関係を扱う点で、
これらの研究と接点をもつ。
ただし本論文は、規範がエージェントに内在することを主張せず、
可行性に関わる構造的制約を
自己維持（SM）、環境利用（EU）、可行応答（FR）の
三原理として分けて定式化する。
また、Maturana と Varela は
生命系の自律性を自己産出（オートポイエーシス）の概念で論じた
\cite{MaturanaVarela1980}。
本論文の構造的自己（定義~\ref{def:s_struct}）は
7タプルとして定め、自己産出を条件に含めない。
Montévil と Mossio は、
生物の組織を、過程に作用する制約が互いに依存し合い互いを維持する閉包として特徴づけ、
これを相互作用する生物系の境界を定める手がかりとしても用いた
\cite{MontevilMossio2015}。
本論文の崩壊の条件 (SC-1)（定義~\ref{def:collapse}）は、
構造的同一性を成立させる組織関係の喪失を述べるが、
その組織関係を制約の閉包として特定することは定義に含めない。

自己と環境の境界、内部状態、行動、存続を一つの説明の体系で扱う研究として、
自由エネルギー原理がある。
Friston は、マルコフブランケットによって内部状態と外部状態が分けられる確率的な力学系では、
内部状態が自由エネルギーと呼ばれる量を最小化するように見え、
環境をモデル化してそれに作用することで自己の構造と機能を保つ
能動的な推論を行っているように見えると論じた
\cite{Friston2013}。
本論文も、自己--環境境界（定義~\ref{def:b_se}）、
内部モデルとその更新（定義~\ref{def:internal_model}、定義~\ref{def:update_operator}）、
行動、および存続を一つの形式で扱う。
ただし本論文の形式は決定論的であり、
時刻別の可行性を持続評価関数の符号条件 $\Pi(x_t)\ge0$ によって表し、
何らかの量の最小化を原理として置かない。
また、Kolchinsky と Wolpert は、
系が環境について持つ情報のうち、
系の存続の度合いを表す可行性関数の値を保つのに因果的に必要なものを意味情報と定義した
\cite{KolchinskyWolpert2018}。
本論文の必要環境要因（定義~\ref{def:environment_space}）は、
存続に必要な外部環境の条件を、
状態と外部環境状態の組に対する充足判定写像として与えるものであり、
系が環境について保持する情報を定義しない。
Ruffini は、決定論的で閉じた世界と調節器の組について、
アルゴリズム的な複雑さによる調節の性能指標と、
調節器が世界について持つ情報との関係を分析した
\cite{Ruffini2026}。
本論文は内部モデル $M$ を7タプルの成分として置くが（定義~\ref{def:internal_model}）、
この性能指標を採用せず、内部モデルの必要性を調節の性能から導かない。
López-Díaz と Gershenson は、
生物の行為者性が、異なる時間尺度で減衰し再生される制約に依存するとして、
自律性、目標指向性、行為者性、開放性の四つの条件を構造的な階層として整理した
\cite{LopezDiazGershenson2026}。
本論文の Phase~1 は、決定論的な離散時刻の標準形に限られ、
このような複数の時間尺度にわたる組織の動学を定式化しない。

本論文の機構記述のうち、
内部の予測と後の時刻の量との照合は、
予測符号化と接点をもつ。
Rao と Ballard のモデルでは、
上位の領野から下位の領野への結合が下位の活動の予測を運び、
下位から上位への結合が予測と実際の活動との残差を運ぶ
\cite{RaoBallard1999}。
\S\ref{sec:tipa} の M1-c は、
予測整合性誤差（定義~\ref{def:m1_temporal_features}）と、
保持記録を通じた予測外の判定（定義~\ref{def:prediction_record_check}）
を記述する。
本論文は、この分野からの機構的導出や
一対一対応を主張するものではなく、
また実証的検証を与えるものでもない。
これらの既存理論との関係を
「ACS理論から導出される」
あるいは
「ACS理論と数学的に同値である」
とは主張しない。
関連分野との対応関係の精密化と検証は、
後続の段階の課題とする。

自動的処理と制御的処理の区別は
Schneider と Shiffrin によって論じられた
\cite{SchneiderShiffrin1977}。
Logan の事例理論は、
格納された事例の想起に依拠する処理を自動的な処理とみなす
\cite{Logan1988}。
これらは、\S\ref{sec:tipa} の Automaticity
（備考~\ref{rem:automaticity_deliberation}）と、
その時刻間の形成・変容
（\S\ref{sec:extended_dynamics} の備考~\ref{rem:s5_fast_change}）
に近い対象を扱う。
本論文は Automaticity を独立した機構として導入せず、
自動的作動と熟慮的作動の間に普遍的な二値境界を置かない。

内部機構の共通の形式を与える試みとして、認知アーキテクチャの研究がある。
Laird、Lebiere、Rosenbloom は、
ACT-R、Sigma、Soar の三つの認知アーキテクチャを踏まえ、
人間に類する心の構造と処理、記憶と内容、学習、知覚と運動について、
共通の計算的な枠組みを提案した
\cite{LairdLebiereRosenbloom2017}。
\S\ref{sec:tipa} の TIPA・PSL・ISL も内部機構の標準形式を与えるが、
可行性に関わる評価を機構の記述に含め、
構成的必要機能（分類A）、標準機能形式（分類B）、記述・診断・比較概念（分類C）を区別する
（定義~\ref{def:abc_classification}）点で目的が異なる。
本論文は、認知アーキテクチャとしての網羅性や実装上の妥当性を主張しない。
Zhang らは、自律的なエージェントを、
定められた同一性と権限の範囲の中で課題の目標と自らの可行性の条件をともに保ち、
世界と自己のモデルを用いて継続的に行動し、
フィードバックを通じて制御の仕組みと長期的な能力を調整する系として定め、
一つの境界と、作動・可行性・反省・学習の四つのループからなる構成を提案した
\cite{ZhangEtAl2026}。
本論文は、TIPA・PSL・ISL とこれらのループとの間に一対一の対応を置かない。
また本論文は、エージェントの構成法を与えるものではなく、
可行性、内部機構、および閉塞と崩壊の区別を、定義と命題として記述する。
Ruffini と Castaldo は、
モデル化・目的関数・計画の三つの機能的役割を備え、
外乱のもとで自らのパターンを保つ調節をもつ系として、
エージェントをアルゴリズム情報理論と熱力学の観点から論じた
\cite{RuffiniCastaldo2026}。
本論文は、TIPA の三つの機構をこれらの役割と一対一に対応させない。

Brooks の階層型の制御では、
下位の層は上位の層が加わっても機能し続け、
上位の層は下位の層の出力を抑制してその役割を包摂しうる
\cite{Brooks1986}。
\S\ref{sec:tipa} の局所保護応答も、
TIPA の出力を用いずに生成される
（備考~\ref{rem:psl4_concurrency}）。
ただし、局所保護応答の発動は M3 の行動指令によって取り消されず、
その局所応答成分は実効行動に保持される
（定義~\ref{def:effective_action_resolve}）。
この点で、上位の層による出力の抑制を許す構成とは異なる。

強化学習の安全性の研究では、
仕様から合成した遮蔽機構（shield）が、
学習器に安全な行動の一覧を与える構成と、
学習器が選んだ行動が仕様に反する場合にだけそれを置き換える構成とが
提案されている \cite{AlshiekhEtAl2018}。
\S\ref{sec:tipa} の実効行動解決は、
TIPA の行動指令とは別の経路で生じる局所応答成分を実効行動に保持する点で
後者と比べられる。
ただし本論文は、仕様の合成や安全性を確保する手法を与えない。

De Jaegher と Di Paolo は、
相互作用の過程そのものが一種の自律性をとりうると論じた
\cite{DeJaegherDiPaolo2007}。
これに対し、\S\ref{sec:rel_min} の双方向合成
（定義~\ref{def:bidirect_relation}）は二つの単方向記録の順序対であり、
相互作用を新しいACSとして構成しない。
複数のACSが参加する組織を一つの系として扱う場合は、
\S\ref{sec:org_acs} の組織水準のACS構造の問題となる。

Castelfranchi は、
他の行為者の心についての信念に基づく弱い社会的行為と、
目標に基づく強い社会的行為を区別し、
目標の委任と採用を協力の構成要素として分析した
\cite{Castelfranchi1998}。
\S\ref{sec:rel_min} の関係カテゴリー（定義~\ref{def:category_set}）も
依存・協力・競争などの名称を含むが、
これらは関係形式を記述する名称であり、
成立条件として行為者の信念や目標の帰属を定めない。
また、関係品質は記述者が定める評価規則に相対する量であり（定義~\ref{def:Q_rule}）、
関係の継続中の可行性関連量の差分を、
その外部対象による因果効果とはしない。

Beer は、組織を管理サイバネティクスの立場から論じた
\cite{Beer1972}。
\S\ref{sec:org_acs} は、
組織にACSが参加することと組織そのものがACSであることを区別し、
組織水準の構造について先行節の定義を適用できる条件を記述する。
特定の統治形態や管理の方式は扱わない。
Masumori は、過程が互いの存続に寄与する関係を、
系の可行性への介入による効果によって定め、
自律性、行為者性、および自己を維持する単位どうしの整合（alignment）を
この関係によって論じた
\cite{Masumori2026}。
本論文の関係記述は、そのような介入による寄与を定義するものではなく、
関係の継続中の可行性関連量の差分を因果効果としない（\S\ref{sec:rel_min}）。
協調的な行動が有効な選択肢として残る条件と、
有効な選択肢が協調的な行動に絞られる条件の分析は、
問題~\ref{prob:capability_paradox} および問題~\ref{prob:ecit_full} に残される。

外乱を伴う系の可行性については、
一方の主体が他方の主体のすべての作用に対して状態を制約集合に保てる
初期状態の最大集合が差別化核（discriminating kernel）として論じられ、
その離散時間版と近似の算法が与えられている
\cite{CardaliaguetQuincampoixSaintPierre1994}。
定義~\ref{def:viab_core} の条件 $(\ast)$ は、
これと同じ $\exists u\,\forall e$ の型の量化条件を用いるものであり、
本論文はこの型の量化を新たに導入するものではない。
ただし、状態空間、外乱の範囲の定め方、および戦略や情報の構造の前提は
文献ごとに異なりうるので、本論文は差別化核との定義の同一性を主張しない。
本論文の構成上の特徴は、外乱の範囲を状態ごとの必要環境充足集合
$E_{\mathrm{req}}(x)$（定義~\ref{def:E_req}）に取り、
その範囲を必要環境要因 $\eta\in E$ と構成要件 CR-1 によって定める点にある。
外部環境状態空間 $E_{\mathrm{ext}}$ のすべての状態に対する頑健な核、
および楽観的な核との包含関係は命題~\ref{prop:kernel_standard_comparison} に示す。
Yamaki と Churiki は、
有限の格子環境で有限の地平の方策を評価する計算モデルを用いて、
生物の認知に類する方策選択を、
可行性の制約のもとで制御可能な将来を広げることとして捉えうるかを検討した
\cite{YamakiChuriki2026}。
本論文の三原理は方策を選ぶ目的関数ではなく（\S\ref{sec:three_principles}）、
FR の時刻別充足述語は将来の制御可能性の大きさを条件に含めない（定義~\ref{def:FR}）。
本論文は、制御可能な将来の拡大から、
可行性の維持を一般的な帰結として置かない。
また、安全制御の分野では、安全集合を関数の上位集合として与え、
その前方不変性を制御によって保つ制御障壁関数の枠組みが用いられている
\cite{AmesEtAl2019}。
可行集合 $K$ と可行行動集合 $A_t$ はこれと同じ形の集合を
一段階の可行性について扱うが、
本論文は安全集合の設計法や制御則の構成を与えない。

可行性理論は回復力の定義にも用いられている。
Martin は、撹乱を受けた後の状態から所与の時間地平までに制約集合へ戻るための
最小の費用を回復の費用とし、その逆数によって回復力を定めた
\cite{Martin2004}。
そこでは、制約を保ち続けられる状態、
所与の時間地平までに有限の費用で回復できる状態、
および回復できない状態が、
回復の費用がゼロ、有限、無限であることによって区別される。
本論文の閉塞（定義~\ref{def:blockage_sec2}）は、
実現した外部環境状態のもとで一段階の可行行動の集合が空であることであり、
崩壊（定義~\ref{def:collapse}）は、
構造的同一性を成立させる組織関係の喪失と、
実現した外生背景成分の列のもとでの通常の自律的更新と
行動応答成分の遷移 $g_O$ による継続、および
許容外的作用（定義~\ref{def:admissible_external_action}）を加えた別の継続に相対する
同一個体としての非回復性によって定まる。
本論文は、回復の費用が無限であることと崩壊とを同一視しない。

可行性の概念を人工的なエージェントに用いる研究もある。
Tsiokos は、資源台帳による実行可能性の制限、頑健な可行性核、
通信路容量としての実行可能なエンパワーメント、経験的なパッケージング自己写像の
四つの検査可能な構成要素によって行為者を定め、
その可行性核を後続台の意味論のもとでの最大不動点として計算した
\cite{Tsiokos2026}。
Marin と Chaudhary は、Aubin の可行性理論に基づいて、
自律的なAIエージェントの実行時の統治の枠組みを提案した
\cite{MarinChaudhary2026}。
本論文の $E_{\mathrm{req}}$ ロバスト可行性核
（定義~\ref{def:viab_core}）は、
$E_{\mathrm{req}}(x)$ のすべての環境に対して自身にとどまる行動が存在する
最大の部分集合として定め、これらの可行性核との同一性を主張しない。
また本論文は、行為者性の指標や統治の手続きを与えるものではなく、
ACSの内部機構と、閉塞・整形式実現の不成立・崩壊の区別
（定義~\ref{def:blockage_sec2}、定義~\ref{def:wf_realization}、
定義~\ref{def:collapse}）を記述する。

Dalrymple らは、AIが外界に及ぼす影響のモデル、
許容される結果の形式的な仕様、
仕様が満たされることの監査可能な証拠を与える検証の三つによって、
AIシステムの安全性について定量的な裏づけを得ることを目指す枠組みを提案した
\cite{Dalrymple2024}。
本論文はAIの安全性を確保する手法を与えるものではなく、
本論文の記述語彙をこうした仕様の記述に用いることは
本論文の主張に含めない。

以上の関連理論は、
本論文の全体を単独で包含する競合理論としてではなく、
可行性、非平滑境界、
エージェント行動、学習、自己維持、
自律性と調節、組織の閉包、自由エネルギー原理と意味情報、
予測、自動的処理、認知アーキテクチャ、階層型の制御、相互作用、社会的行為、組織、
外乱のもとでの可行性核、安全集合と行動の遮蔽、可行性に基づく回復力、恒常性に基づく強化学習、
AIエージェントの可行性と安全性等の
個別側面について、接続する既存基盤または比較の対象として参照される。
本論文の目的は、
これらの既存理論を置き換えることではなく、
それぞれが扱う一部の概念との接点・相違点を明示しつつ、
ACSの可行性、内部機構、時刻間動学、
関係性、失敗状態を
一つの決定論的形式枠組みへ接続することにある。

%----------------------------------------------------------------------
\subsection{本論文の構成}
\label{subsec:roadmap}

本論文の構成は次のとおりである。

\begin{itemize}

  \item \S\ref{sec:acs_def} では、
    ACSを7タプルとして形式的に定義し、
    持続評価関数、可行集合、可行行動集合、
    内部モデル、更新作用素、認知的処理能力の記述方針、
    構造的自己その他の基本語彙を確立する。

  \item \S\ref{sec:three_principles} では、
    三原理SM・EU・FRの時刻別充足述語を定義し、
    それぞれの論理的位置づけ、
    閉塞の二分、時刻別充足述語の間の条件付き含意と含意関係の限定、
    応答多様性縮退と長期可行性の条件付き接続、
    および三原理を有限時間窓上で統合する
    三原理時間統合評価条件（T-STAB）を扱う。

  \item \S\ref{sec:tipa} では、
    原理水準の三原理を再定義せず、
    本論文が対象とするACSに必要な機能を
    分類Aとして必要に応じて明示し、
    それらを本論文内で分析可能にする
    分類Bの標準機能形式として、
    TIPA、PSL、ISLからなる内部機構構造を定式化する。
    あわせて、識別、評価、応答形成、
    機能的記憶、高速処理、Automaticity、
    作動水準等の責務と境界を整理する。

  \item \S\ref{sec:extended_dynamics} では、
    \S\ref{sec:tipa} の同一時刻構造を再定義せず、
    一段の作動記録が次の時刻へ閉じる標準作動形式を与え、
    単一ACSにおける内部モデル、
    機能的記憶、各種記述量、
    PSL・ISLの健全度等の
    決定論的な時刻間変化を整理する。
    また、分類Cの記述・診断・比較量が
    ACSの追加構成要素ではないことを明示する。

  \item \S\ref{sec:rel_min} では、
    ACSが外部対象と取り結ぶ関係性の最小構造を、
    関係継続中に記録される可行性関連量の変化を記述する形式として扱う。
    外部対象の分類、単方向・双方向の関係記述、
    評価規則に相対する関係品質、
    関係カテゴリーおよび方向性を、
    Phase~1 の範囲で価値中立的に整理する。

  \item \S\ref{sec:collapse_cond} では、
    \S\ref{sec:acs_def} で定義した閉塞と崩壊を区別し、
    可行行動集合の消失と
    構造的自己の回復不能性を
    異なる失敗状態として扱う。
    また、外的回復可能性と、それらへの到達経路を
    Phase~1 の決定論的枠組みの範囲で整理する。

  \item \S\ref{sec:org_acs} では、
    単一ACSの枠組みを
    単なるACS集団と区別された組織ACSへ適用し、
    組織水準のACS構造の成立のもとで
    先行節の定義の組織水準への適用を扱う。

  \item 最後に \S\ref{sec:limitations_and_outlook} で、
    Phase~1 の理論的限界、Phase~1.$\alpha$ への委譲事項、未解決問題を開示し、
    Phase~2・Phase~3への展望と結論を述べる。

  \item 付録~\ref{app:notation} に記号添字規約を、
    付録~\ref{app:symbols} に各節で新たに用いる主要記号の一覧を、
    付録~\ref{app:index} に定義・仮定・命題等の一覧を置く。

\end{itemize}
